% ==============================================================
% METHODOLOGY EXTENSION: Highway Upgrades and the Braess Effect
% Insert this as Section 3.X after Section 3.1 
% (Network Representation and Time-Dependent Travel Times)
% and update Section 3.2 (Multi-Objective Path Formulation)
% accordingly using the revised equations below.
% ==============================================================

% ---------------------------------------------------------------
\section{Highway Upgrades, Structural Vulnerability, and the
         Braess Paradox in Organic Urban Networks}
\label{sec:braess}
% ---------------------------------------------------------------

A distinctive feature of poorly developed cities in sub-Saharan Africa,
including Zomba and Blantyre, is that road investment is incremental and
spatially uneven. Individual corridors are periodically upgraded to trunk
or dual-carriageway standard while the surrounding network remains
unchanged. This pattern produces a structural condition that is
theoretically familiar yet practically undermodelled in the emergency
routing literature: the conditions described by Braess's paradox
\cite{braess1968}.

Braess's paradox demonstrates that adding capacity to a congested network
can, under user-equilibrium routing, increase the total travel time
experienced by all users. The mechanism is not that the new road is
poorly built; it is that its existence induces diversion. Drivers who
previously distributed across several routes now converge on the new
high-capacity corridor, and the feeder roads that channel them into it
absorb the concentrated demand. The highway flows freely; the connectors
that feed it become the new bottleneck. For emergency vehicles, which
cannot deviate from the network in the way a car might, the consequence
is that routes passing through feeder segments adjacent to recently
upgraded corridors become systematically less reliable, not because their
own physical condition has changed, but because the structural load
imposed on them has increased as a downstream effect of an improvement
elsewhere in the network.

This section formalises how this effect is incorporated into the
multi-objective model without adding a fourth objective dimension or
restructuring the MOLS algorithm.

\subsection{Highway Edge Classification and Constant Travel Time}
\label{subsec:hw_classification}

We partition the edge set $E$ into three disjoint classes:

\begin{definition}[Edge Classification]
\label{def:edge_class}
\begin{align}
  E_H &= \{ (i,j) \in E : \text{edge is a trunk/highway segment} \},
                                                         \label{eq:EH}\\
  E_F &= \{ (i,j) \in E\setminus E_H :
           i \in \mathcal{N}_H \text{ or } j \in \mathcal{N}_H \},
                                                         \label{eq:EF}\\
  E_S &= E \setminus (E_H \cup E_F),                    \label{eq:ES}
\end{align}
where $\mathcal{N}_H = \bigl\{ v \in V : \exists\,(i,j)\in E_H,\;
i=v \text{ or } j=v \bigr\}$ is the set of nodes directly incident to
at least one highway edge.
\end{definition}

Edges in $E_H$ represent either existing trunk roads or primary roads
that have been upgraded to highway standard. Edges in $E_F$ are
\textit{feeder} segments: roads that are not highways themselves but
whose endpoints connect directly to the highway corridor, and therefore
absorb the redistribution of traffic that the paradox predicts. All
remaining edges belong to $E_S$.

\textbf{Highway travel time.}~
Unlike ordinary road segments, trunk-standard corridors in sub-Saharan
cities experience congestion only at extreme traffic volumes that are
rarely reached in secondary cities \cite{gonzalez2023}. We therefore
model highway edges with time-invariant travel time:

\begin{equation}
  \tau_{ij}(t) = \tau_{ij}^{0}
  = \frac{\ell_{ij}}{s_{ij}^{0}}, \qquad \forall\,(i,j)\in E_H,
  \quad \forall\, t,
  \label{eq:hw_const}
\end{equation}

where $\ell_{ij}$ is the edge length and $s_{ij}^{0}$ is the free-flow
speed. This is a deliberate modelling choice, not a universal claim about
highway behaviour: it reflects the low traffic volumes typical of
secondary African cities and ensures that the paradox effect is isolated
to the feeder segments where it actually manifests, rather than being
diffused across the whole network.

\begin{remark}
Equation~\eqref{eq:hw_const} is implemented by assigning a
congestion multiplier of $1.0$ across all time periods for $E_H$ edges,
leaving the piecewise-linear structure of Section~\ref{sec:methodology}
intact. No algorithmic change to the MOLS procedure is required.
\end{remark}

\subsection{The Braess Structural Vulnerability Parameter}
\label{subsec:phi}

For feeder edges $(i,j)\in E_F$, the standard variability measure
$v_{ij}(t_i)$ in objective $f_2$ underestimates true unreliability
because it captures only the effect of current-period congestion. It does
not capture the structurally elevated load that results from being a
collector for a high-capacity corridor. We introduce a time-independent
\textit{Braess structural vulnerability parameter} $\varphi_{ij}\in
[0,1]$ to represent this additional source of unreliability.

\begin{definition}[Braess Structural Vulnerability Parameter]
\label{def:phi}
For each edge $(i,j)\in E$:
\[
  \varphi_{ij} =
  \begin{cases}
    0, & (i,j)\in E_S\cup E_H, \\[4pt]
    \varphi_{ij}^{*}\in(0,1], & (i,j)\in E_F,
  \end{cases}
\]
where $\varphi_{ij}^{*}$ is the edge-specific vulnerability score,
calibrated by the proximity of the feeder to the highway node and
the proportion of network shortest paths that pass through both
the feeder and the highway \cite{6,11}.
\end{definition}

The parameter is time-independent because it reflects a structural
property of the network, not a time-of-day pattern. A road that feeds a
highway is permanently more exposed to load redistribution once the
highway is present; this does not switch off at night.

\subsection{Modified Reliability Objective}
\label{subsec:f2_modified}

The reliability objective $f_2$ is updated to incorporate
$\varphi_{ij}$:

\begin{equation}
  \boxed{
  f_2(P) = \sum_{(i,j)\in P} v_{ij}^{*}(t_i),
  \quad
  v_{ij}^{*}(t_i)
    = \tau_{ij}(t_i)
      \cdot\!\left(\frac{\tau_{ij}(t_i)}{\tau_{ij}^{0}}\right)^{\!2}
      \cdot \gamma_{ij}
      \cdot \bigl(1 + \varphi_{ij}\bigr).
  }
  \label{eq:f2_modified}
\end{equation}

The multiplicative factor $(1+\varphi_{ij})$ inflates the variability
penalty for feeder edges without changing its functional form, so that
the MOLS label update remains:

\begin{equation}
  \hat{f}_2 \;\leftarrow\; f_2
    + \tau_{ij}(t)\cdot\!\left(\frac{\tau_{ij}(t)}{\tau_{ij}^{0}}
      \right)^{\!2}\!\cdot\gamma_{ij}\cdot(1+\varphi_{ij}),
  \label{eq:mols_f2_update}
\end{equation}

requiring only a single additional lookup per edge extension in
Algorithm~\ref{alg:mols}, with no change to the dominance criterion or
label structure. The complete three-objective formulation now reads:

\begin{equation}
  \begin{aligned}
    f_1'(P) &= \sum_{(i,j)\in P} \bigl[\tau_{ij}(t_i) + \lambda\,
               c_{ij}\bigr],
               && \text{(travel time + cost)} \\[4pt]
    f_2(P)  &= \sum_{(i,j)\in P}
               \tau_{ij}(t_i)
               \!\left(\frac{\tau_{ij}(t_i)}{\tau_{ij}^{0}}\right)^{\!2}
               \!\gamma_{ij}(1+\varphi_{ij}),
               && \text{(reliability + Braess penalty)} \\[4pt]
    f_3(P)  &= \beta \sum_{(i,j)\in P} \rho_{ij},
               && \text{(resilience)}
  \end{aligned}
  \label{eq:mop_revised}
\end{equation}

\begin{remark}
The Braess penalty operates through $f_2$ rather than as an independent
fourth objective for two reasons. First, it is a source of travel time
unreliability, not a conceptually separate concern; consolidating it with
$\gamma$ preserves the interpretability of the Pareto front. Second,
adding a fourth objective dimension would increase the size of the
non-dominated set combinatorially, imposing computational costs that are
unjustified given the binary character of feeder classification and the
single calibration parameter involved. A fourth objective formulation
would be appropriate if $\varphi_{ij}$ were derived from a continuous
field measurement and carried distinct decision-making weight, which is
not the case here.
\end{remark}

\begin{remark}
The interaction between $f_1'$ and $f_2$ under this formulation is
behaviorally meaningful. Highway edges $(E_H)$ reduce $f_1'$ by
providing fast constant-time travel, while feeder edges $(E_F)$ inflate
$f_2$ through $(1+\varphi_{ij})$. The MOLS Pareto front therefore
naturally surfaces the Braess trade-off: routes that use the highway are
fast but their reliability cost is higher because they must traverse
feeder segments to access it. Decision-makers operating under tight time
constraints (critical emergencies, $\beta$ small, reliability weight low)
will rationally select routes exploiting the highway; those prioritising
reliability under chronic congestion will tend to avoid feeder segments
even at the cost of longer nominal travel time.
\end{remark}

\subsection{Calibration of $\varphi_{ij}$}
\label{subsec:phi_calibration}

In this study $\varphi_{ij}$ values for feeder edges are assigned
synthetically. We define the vulnerability of a feeder edge as
proportional to the betweenness centrality of the highway node it
connects to, normalised within the feeder set:

\begin{equation}
  \varphi_{ij}^{*}
  = \varphi_{\min}
    + (\varphi_{\max}-\varphi_{\min})
      \cdot \widetilde{\rho}_{\mathcal{N}_H(ij)},
  \label{eq:phi_calibration}
\end{equation}

where $\widetilde{\rho}_{\mathcal{N}_H(ij)}$ is the normalised
betweenness centrality of the highway-incident node connected to edge
$(i,j)$, and we set $\varphi_{\min}=0.25$, $\varphi_{\max}=0.70$.
These bounds reflect the empirical range of load-redistribution effects
reported in the Braess-paradox literature for networks of comparable
size and density \cite{11}.

For real-world deployment, $\varphi_{ij}$ can be estimated from GPS
speed data collected before and after a highway-upgrade event, using
the ratio of post-upgrade to pre-upgrade travel time variance on each
feeder link as a direct empirical proxy. In the absence of such data,
the structural approximation in \eqref{eq:phi_calibration} provides a
principled lower bound on feeder vulnerability.

% ---------------------------------------------------------------
% UPDATE TO PARAMETERS TABLE (replace or extend Table~\ref{tab:params})
% ---------------------------------------------------------------
\begin{table}[H]
  \centering
  \caption{Extended Model Parameters Including Braess Vulnerability}
  \label{tab:params_ext}
  \begin{tabular}{llll}
    \hline
    \textbf{Parameter} & \textbf{Symbol} & \textbf{Value}
                       & \textbf{Justification} \\
    \hline
    Resilience weight        & $\beta$           & 0.5
      & Balanced time/resilience trade-off \cite{6,11} \\
    Cost conversion factor   & $\lambda$         & 0.001
      & MWK to minute-equivalent \cite{tirkolaee2020} \\
    Variability factor       & $\gamma$          & 0.25--0.60
      & Per-edge (reliability\_params.csv) \\
    Congestion update        & $\Delta$          & 15\%
      & Vehicle interaction increment \cite{luan2024} \\
    Market peak multiplier   & ---               & 2.3--2.5$\times$
      & Morning/evening peaks \cite{8,9} \\
    Normal peak multiplier   & ---               & 1.5$\times$
      & Morning/evening peaks \cite{13} \\
    Criticality threshold    & $\rho_0$          & 0.3
      & High-betweenness classification \cite{11} \\
    Braess penalty (feeder)  & $\varphi_{ij}$    & 0.25--0.70
      & Feeder vulnerability; 0 for $E_S\cup E_H$
        (\S\ref{subsec:phi_calibration}) \\
    Highway const.\ tt.      & $\tau_{ij}^0$     & free-flow only
      & Eq.~\eqref{eq:hw_const}; 6 edges in $E_H$ \\
    \hline
  \end{tabular}
\end{table}

% ---------------------------------------------------------------
% REQUIRED ADDITION TO BIBLIOGRAPHY
% ---------------------------------------------------------------
% Add this entry to your \begin{thebibliography} block:
%
% \bibitem{braess1968}
% Braess, D. (1968). \"Uber ein Paradoxon aus der Verkehrsplanung.
% \textit{Unternehmensforschung}, 12(1), 258--268.
% [English translation: Braess, D., Nagurney, A., \& Wakolbinger, T.
% (2005). On a paradox of traffic planning.
% \textit{Transportation Science}, 39(4), 446--450.
% \url{https://doi.org/10.1287/trsc.1050.0127}]
